% §17.7 Bridge
%
% Cross-reference resolution status:
%
% DEFINED here:
%   \label{sec:ch17-bridge}                   -- this section
%
% INTERNAL refs (chapter 17):
%   \Cref{sec:what-is-cipher-map}             -- §17.1
%   \Cref{sec:bernoulli-error-model}          -- §17.2
%   \Cref{sec:trapdoor-paradigm}              -- §17.3
%   \Cref{sec:fhe-sidebar}                    -- §17.4
%   \Cref{sec:boolean-trapdoor}               -- §17.5
%   \Cref{sec:maximizing-confidentiality}     -- §17.6
%   \Cref{sec:bib-notes-ch17}                 -- §17.8 (Task 9)
%
% RESOLVED (chapter labels for forward refs):
%   \Cref{ch:cipher-maps-trapdoor}            -- this chapter
%   Chapter 18 (encrypted search)             -- Plan 0 stub
%
% §17.7 closes chapter 17 with three items: (1) recap of the dual
% role of fprate/fnrate (correctness AND confidentiality);
% (2) chapter 18 forward-link to encrypted search as a specific
% cipher-map family; (3) signpost to §17.8 antecedents.

\section{Bridge}
\label{sec:ch17-bridge}

The chapter has built a single bridge in six pieces. The cipher-map
abstraction of \Cref{sec:what-is-cipher-map} reframed the
encryption setting as a pair of error rates rather than a pair of
keys: the server's job is to approximate $f$ on encrypted inputs,
and the bounded $\fprate$ and $\fnrate$ control how well it
succeeds. The error-model section
(\Cref{sec:bernoulli-error-model}) identified those rates with the
Bernoulli atom of chapter~4 and lifted chapter~4's Kronecker
composition law to chained cipher-map evaluation. The
trapdoor-paradigm section (\Cref{sec:trapdoor-paradigm}) introduced
the client-server split that makes the abstraction useful in
practice, and the FHE sidebar (\Cref{sec:fhe-sidebar}) marked the
zero-error baseline against which the cipher-map approach is
measured. Sections~\ref{sec:boolean-trapdoor}
and~\ref{sec:maximizing-confidentiality} pushed the framework
forward in two directions: Boolean composition over trapdoor sets
(chapter~3's composition theorem in cryptographic clothing) and
confidentiality as a chapter~7 PPV.

\subsection*{The dual role of the error rates}

The recurring thread is the dual role of $\fprate$ and $\fnrate$.
In Parts~I-IV the two rates were operational parameters: they
described how often the framework's structures returned wrong
answers, and the design problem was to keep them small subject to
storage and runtime budgets. In chapter~17 the same parameters
acquired a second job. The cipher-map false-positive rate is also
the share of the adversary's $\top$-observations that fail to
identify members of $A$, and the cipher-map false-negative rate is
also the share of $\bot$-observations that protect real members.
The two rates together fix the cipher map's PPV at workload
prevalence $\pi$, and the complementary $1 - \PPV$ is the
confidentiality margin.

The duality is not a coincidence. The Bernoulli framework
identifies a structure by its $(\fprate, \fnrate)$ point in the
Bernoulli square; chapter~6 read that point geometrically as a
position on the perfect-filter line; chapter~7 read the same point
predictively as PPV and NPV. Chapter~17 has read the point a third
time, cryptographically, as the trade-off between operational
fidelity and informational opacity. The three readings are the
same algebra. Designing a cipher map means choosing where to live
in the Bernoulli square, and that single choice now answers three
different questions at once.

\subsection*{Forward to chapter~18: encrypted search}

Chapter~18 takes the cipher-map machinery and specializes it to
one application: encrypted search. The setting fixes $f$ to a
membership-test family ``does record $r$ match query $q$?'' and
asks how a server that holds an encrypted index can answer such
queries faithfully without revealing $r$ or $q$ to whoever is
running the server. Bloom-filter cipher maps from
\Cref{sec:boolean-trapdoor} reappear as the simplest index
structure, and the chapter~14 BHF construction reappears as a
space-optimal alternative. The two-party setup of
\Cref{sec:trapdoor-paradigm} is still in force; the
$(\fprate, \fnrate)$ point still does triple duty as
correctness, geometry, and confidentiality.

What chapter~18 adds is the structure that the encrypted-search
literature has built around the membership-test case: oblivious
access patterns, side-channel considerations, multi-query leakage
bounds, and the relation between confidentiality $C$ and the
information-theoretic measures the encrypted-search literature
prefers. Chapter~17's algebraic machinery sets up the formal
language; chapter~18 fills in the application-specific structure.
The conceptual move from chapter~17 to chapter~18 is the same
move chapters~14 and~15 made together: chapter~17 is the abstract
algebra, chapter~18 is the concrete instance.

\subsection*{Antecedents}

The cipher-map abstraction and the trapdoor-computing paradigm
have antecedents in the cryptographic literature that the chapter
has named only in passing. The FHE-sidebar section pointed to
Gentry's 2009 construction as the zero-error baseline; the
trapdoor-paradigm section borrowed terminology from the
oblivious-computing tradition that chapter~16's bibliographic
notes already touched. The next section
(\Cref{sec:bib-notes-ch17}) attaches the chapter's prose to these
antecedents formally and gives the reader entry points to the
four literatures the chapter has been drawing on.
